\chapter{変分 Wasserstein 問題}
\label{c-variational}

データ分析において、確率測度間の一般的なダイバージェンス（\eg ユークリッド、全変動、Hellinger、Kullback--Leibler）は、フィッティング誤差やパラメータ推定問題における損失の測定によく用いられる。本章まで、我々は optimal transport 幾何学が、他の情報ダイバージェンスとは共有されない独自の能力、すなわち物理的なアイデア（質量の変位）と幾何学（観測値やビン間の ground cost）を活用して測度を比較する能力を持つことを論じてきた。これら2つの事実を組み合わせると、Wasserstein 距離を損失関数として使用することは非常に魅力的である。このアイデアは最近、さまざまな応用問題に対して探求されている。しかし、このアイデアに伴う主な技術的課題は、Wasserstein 距離を効率的に近似し微分することにある。本章ではまず、いくつかの動機付けとなる例を示し、本書の前半の章で提示されたさまざまな数値スキームが変分 Wasserstein 問題の解決にどのように使用できるかを示す。

画像処理において、Wasserstein 距離はテクスチャ合成~\citep{2016-tartavel-siims}における損失として用いることができ、合成された統計量と入力例の統計量との間の不一致を考慮する。また、画像領域の統計的均質性を考慮する画像セグメンテーション~\citep{SchnorSegmentation,RabinPapadakisSSVM,peyre2012wasserstein,ni2009local,schmitzer2013object,liimage}にも用いられる。
Wasserstein 距離は、測定値が確率測度である逆問題に対しても非常に自然な忠実度項であり、例えば画像復元~\citep{lellmann2014imaging}、トモグラフィー逆変換~\citep{AbrahamRadon}、密度正則化~\citep{Burger-JKO}、粒子画像速度計測法~\citep{saumier2015optimal}、スパース回復および圧縮センシング~\citep{indyk2011k}、地震波インバージョン~\citep{metivier2016optimal}が挙げられる。
測度間の距離（主に第~\ref{sec-mmd}節で示されるカーネルベースのもの）は、計算解剖学~\citep{vaillant2005surface}における形状マッチング（しばしばカレントと呼ばれるリフトされた空間上の測度として表現）に日常的に用いられるが、OT 距離は興味深い代替手段を提供する~\citep{2017-feydy-miccai}。
ヒストグラムのデータセットの次元削減のために、\citeauthor{lee1999learning}は非負値行列因子分解問題を再構成損失を定量化する Kullback--Leibler ダイバージェンスを用いて定式化できることを示した~\citep{lee1999learning}。ヒストグラムのビンの幾何学に関する事前情報が利用可能な場合、代わりに Wasserstein 距離を用いることができ、著しく異なる結果が得られる~\citep{Sandler09,ZenICPR14,pmlr-v51-rolet16}。

これらの問題のすべてに共通するのは、Wasserstein 距離の勾配またはその近似へのアクセスが必要であることである。本節では、まずそのような勾配を近似する方法を提示し、次に変分 Wasserstein 問題として定式化できる3つの重要な応用を紹介する。

%%%%%%%%%%%%%%%%%%%%%%%%%%%%%%%%%%%%%
\section{Wasserstein 損失の微分}

統計学、テキスト処理、または画像処理において、通常、測定から得られる確率分布 $\be$ をモデル、すなわちパラメータ化された分布の族 $\{\al_\th,\th\in\Theta\}$ と比較する必要がある。ここで $\Theta$ はユークリッド空間の部分集合である。このような比較は「損失」または「忠実度」項を通じて行われ、本節では Wasserstein 距離がその役割を果たす。
最も単純なシナリオでは、適切なパラメータ $\th$ の計算は、以下を直接最小化することで得られる：
\eql{\label{eq-wloss-generic}
	\umin{\th \in \Theta} \Ee(\th) \eqdef \MK_\c(\al_\th,\be).
}
もちろん、より複雑な問題を考えることもできる。例えば、第~\ref{sec-bary}節で述べる重心問題はこのような項の和で構成される。しかし、これらのより高度な問題のほとんどは、上記の基本的な場合に対して定義されたツールを、連鎖律を用いて明示的に導関数を計算するか、注意~\ref{rem-auto-diff}で提唱される自動微分を用いることで、適応させることによって解くことができる。

\paragraph{凸性.} 2つのヒストグラムまたは2つの密度間の Wasserstein 距離は、それぞれ~\eqref{eq-dual}および~\eqref{eq-dual-generic}に示されるように、2つの入力に関して凸である。したがって、パラメータ $\theta$ 自体がヒストグラムである場合、すなわち $\Theta=\simplex_n$ かつ $\al_\theta=\theta$ の場合、またはより一般的に $\theta$ が単体における $K$ 個の重みを記述する場合、$\Theta=\simplex_K$ かつ $\alpha_\theta=\sum_{i=1}^K \theta_i \alpha_i$ が $\simplex_N$ における既知のアトム $\alpha_1,\dots,\alpha_K$ の凸結合である場合、問題~\eqref{eq-wloss-generic}は凸のままである（前者は重心問題に対応し、後者は Wasserstein 損失を用いた辞書学習問題の1回の反復に対応する~\citep{pmlr-v51-rolet16}）。しかし、より一般的なパラメータ化 $\th\mapsto \al_\th$ に対して、問題~\eqref{eq-wloss-generic}は一般に凸ではない。


\paragraph{単純な場合.} Wasserstein 距離が閉じた形を持つ単純な場合、例えば1変量（注意~\ref{rem-1d-ot-generic}を参照）や楕円等高線（第~\ref{rem-dist-gaussians}節を参照）分布に対しては、単純な回避策が存在する。これらは主に、Wasserstein 距離をこれらの分布の適切な表現間のより単純な距離（1変量測度に対する分位関数上のユークリッド距離、同じ族の楕円等高線分布に対する共分散行列上の Bures 計量）として表し、そのような表現上で問題~\eqref{eq-wloss-generic}を直接解くことからなる。

しかし、ほとんどの場合、問題~\eqref{eq-wloss-generic}の局所最小化を計算するために $\al_\th$ の慎重な離散化に頼る必要がある。2つのアプローチが考えられる：Euler 的方法と Lagrange 的方法である。図~\ref{fig-eulerian-lagrangian}はこれら2つの基本的な離散化スキームの違いを示している。この議論を過度に単純化する危険を冒して言えば、Euler 的離散化は測度が低次元空間上に台を持つ場合（形状や色空間を扱う場合など）、または本質的に離散的な問題（文字列やテキスト分析から生じるものなど）に最も適している。観測が高次元空間で連続的な値を取りうるフィッティング問題に適用する場合、Lagrange 的な視点が通常、唯一の適切な選択肢である。


\begin{figure}[h!]
\centering
\includegraphics[width=\linewidth]{eulerian-lagrangian/eulerian-lagrangian.pdf}
\caption{\label{fig-eulerian-lagrangian}
密度を持つ連続分布（紫、左）の、Lagrange 的表現 $\frac{1}{n}\sum_i \de_{x_i}$（青、上）とサイズが増加する格子上のセルを表す $x_i$ を用いた Euler 的表現 $\sum_i \a_i \de_{x_i}$（赤、下）による、次第に細かくなる離散化。Euler 的な視点はピクセル化された画像から始まり、元の密度にほぼ一致するほどの細かい解像度へと進む。重み $\a_i$ は各ピクセルセルの強度に直接比例する。
}
\end{figure}


%%%%%%%%%%%%%%%%%%%%%%%%%%%%%%%%%%%%%
\subsection{Euler 的離散化}

問題を離散化する第1の方法は、両方の分布 $\be=\sum_{j=1}^m \b_j \de_{y_j}$ および $\al_\th=\sum_{i=1}^n \a(\th)_i \de_{x_i}$ が固定された位置 $(x_i)_i$ と $(y_j)_j$ 上に定義された離散分布であると仮定することである。このような位置は、観測空間全体を格子で分割するセルを表すか、連続空間における有限個の関心のある点の集合（例えば、与えられた辞書内の全単語に対するベクトル埋め込みの族~\citep{kusner2015word,pmlr-v51-rolet16}）を表す場合がある。パラメータ化された測度 $\al_\th$ はこの場合、重みベクトル $\a : \th \mapsto \a(\th) \in \Si_n$ によって完全に表現され、格子が大きい場合、実際にはこれは非常にスパースになり得る。この設定は、いわゆる Euler 的離散化手法のクラスに対応する。元の形式では、問題~\eqref{eq-wloss-generic}の目的関数は微分可能でない。滑らかな最小化問題を得るために、エントロピー正則化 OT を用い、~\eqref{eq-wloss-generic}を以下で近似する：
\eq{
	\umin{\th \in \Theta} \Ee_E(\th) \eqdef \MKD_\C^\epsilon(\a(\th),\b)
	\qwhereq
	\C_{i,j} \eqdef c(x_i,y_j).
}
命題~\ref{prop-convexity-dual}がエントロピー損失関数は入力ヒストグラムに関して微分可能かつ凸であり、勾配を持つことを示していることを想起する。

%%%%%%%%%%
\begin{proposition}[ヒストグラムに関する微分]
$\epsilon>0$ に対して、
$(\a,\b) \mapsto \MKD_\C^\epsilon(\a,\b)$ は凸かつ微分可能である。その勾配は
\eql{\label{eq-diff-marginals}
		\nabla \MKD_\C^\epsilon(\a,\b) = (\fD,\gD),
}
で与えられる。ここで $(\fD,\gD)$ は~\eqref{eq-dual-formulation}の一意解であり、$\sum_i \fD_i=\sum_j \gD_j=0$ となるように中心化されている。$\epsilon=0$ の場合、この公式は $\MKD_\C^\epsilon$ の劣微分の要素を定義し、それらが一意であれば関数は微分可能である。
\end{proposition}

$(\fD,\gD)$ の零平均条件は、質量保存を保証するために勾配降下法を用いる際に重要である。
連鎖律を用いると、$\Ee_E$ が滑らかであり、その勾配が
\eql{\label{eq-diff-loss-eul}
	\nabla \Ee_E(\th) = [ \partial \a(\th) ]^\top( \fD\, ),
}
であることがわかる。ここで $\partial \a(\th) \in \RR^{n \times \text{dim}(\Theta)}$ は写像 $\a(\th)$ の Jacobi 行列（微分）であり、$\fD \in \RR^n$ はコスト行列 $\C$ に対する $\a(\th)$ と $\b$ の間の双対エントロピー OT~\eqref{eq-dual-formulation}に付随する双対ポテンシャルベクトルである（Euler 的設定では $\C$ は固定であり、特に $\th$ に依存しない）。
この結果を用いて、勾配降下法により $\Ee_E$ を局所的に最小化できる。

%%%%%%%%%%%%%%%%%%%%%%%%%%%%%%%%%%%%%
\subsection{Lagrange 的離散化}

異なるアプローチは、代わりに固定された（典型的には一様な）重みを用い、入力測度 $\alpha$ を点群パラメータ化写像 $x : \th \mapsto x(\th) = (x(\th)_i)_{i=1}^n \in \Xx^n$ を用いた経験測度 $\al_\th = \frac{1}{n}\sum_i \de_{x(\th)_i}$ として近似することからなる。ここで $\X$ はユークリッド空間であると仮定する。問題~\eqref{eq-wloss-generic}は以下のように近似される：
\eql{\label{eq-lagr-sinkloss}
	\umin{\th} \Ee_L(\th) \eqdef \MKD_{\C(x(\th))}^\epsilon(\ones_n/n,\b)
	\qwhereq
	\C(x)_{i,j} \eqdef c(x(\theta)_i,y_j).
}
ここで、$\al_\th$ の台が $\theta$ とともに変化するため、コスト行列 $\C(x(\th))$ が $\th$ に依存することに注意する。
以下の命題は、エントロピー OT 損失がコスト行列の滑らかな関数であることを示し、その勾配の表現を与える。

%%%%%%%%%%
\begin{proposition}[コストに関する微分]
固定された入力ヒストグラム $(\a,\b)$ に対して、$\epsilon>0$ のとき、写像 $\C \mapsto \Rr(\C) \eqdef \MKD_\C^\epsilon(\a,\b)$ は凹かつ滑らかであり、
\eql{\label{eq-diff-cost}
	\nabla \Rr(\C) = \P,
}
が成り立つ。ここで $\P$ は~\eqref{eq-regularized-discr}の一意の最適解である。
$\epsilon=0$ の場合、この公式は上勾配の集合を定義する。
\end{proposition}
%%%%%%%%%%


$(\X,\Y)$ が $\RR^\dim$ の凸部分集合であり、離散測度 $(\al,\be)$ が~\eqref{eq-pair-discr}の形であると仮定すると、連鎖律を用いて $x = (x_i)_{i=1}^n  \in \X^n \mapsto \Ff(x) \eqdef  \MKD_{\C(x)}(\ones_n/n,\b)$ が滑らかであり、
\eql{\label{eq-diff-positions}
	\nabla \Ff(x) = \pa{ \sum_{j=1}^m \P_{i,j} \nabla_1 \c(x_i,y_j) }_{i=1}^n \in \X^n,
}
であることがわかる。ここで $\nabla_1 \c$ は第1変数に関する勾配である。
例えば、$\X=\Y=\RR^d$ 上で $\c(s,t) = \norm{s-t}^2$ の場合、
\eql{\label{eq-diff-positions-eucl}
	\nabla \Ff(x) = 2 \pa{ \a_i x_i - \sum_{j=1}^m \P_{i,j} y_j }_{i=1}^n,
}
となる。ここで $\a_i=1/n$ である。
定数の差を除き、この勾配は $\Id-\T$ であり、$\T$ は~\eqref{eq-baryproj}で定義される重心射影である。
連鎖律を用いると、Lagrange 的に離散化された問題~\eqref{eq-lagr-sinkloss}が滑らかであり、その勾配が
\eql{\label{eq-diff-loss-lagr}
	\nabla \Ee_L(\th) = [ \partial x(\th) ]^\top( \nabla \Ff( x(\th) ) ),
}
であることがわかる。ここで $\partial x(\th) \in \RR^{\text{dim}(\Theta) \times (n\dim)}$ は写像 $x(\th)$ の Jacobi 行列であり、$\nabla \Ff$ は $\al_\th$ と $\be$ の間の最適カップリング行列を $\P$ として用い、~\eqref{eq-diff-positions}または~\eqref{eq-diff-positions-eucl}のように実装される。
したがって、例えば~\citep{CuturiBarycenter}で用いられているように、勾配降下法を実装して $\Ee_L$ の局所最小化を計算できる。

%%%%%%%%%%%%%%%%%%%%%%%%%%%%%%%%%%%%%
\subsection{自動微分}
\label{rem-auto-diff}

公式~\eqref{eq-diff-loss-eul}および~\eqref{eq-diff-loss-lagr}を適用する際の困難さは、これらの公式が有効であるためには正確な最適解 $\fD$ または $\P$ を計算する必要があることであり、これは非常に多数の Sinkhorn 反復を用いなければ許容可能な精度で達成できない。
比較すべきヒストグラムのサイズと量が大きい困難な状況では、単一の Wasserstein 距離を計算するための計算予算は通常限られており、少数の Sinkhorn 反復しか許容されない。この場合、与えられた反復で得られた値を用いて勾配~\eqref{eq-dual-formulation}を近似するよりも、逆モード自動微分を用いて Sinkhorn アルゴリズムの出力を直接微分する方が通常よい。これは~\eqref{eq-regularized-discr}における正則化 optimal transport のエントロピーを組み込んだ量 $\MKD_\C^\varepsilon$ ではなく、~\eqref{eq-algorithmic-loss}で導入された「アルゴリズム的」Sinkhorn ダイバージェンスを用い、入力（Euler 的場合はヒストグラム、Lagrange 的場合はコスト行列）を用いた単純な写像の合成として直接微分することに対応する。第~\ref{sec-regularized-cost}節で導入された定義を用いると、これはそれぞれ Euler 的および Lagrange 的な場合において、十分大きな $L$ に対して
$$\itL{\SINKHORND_{\C}}(\a(\theta),\b) \quad\text{ or }\quad \itL{\SINKHORND_{\C(x(\theta))}}(\a,\b)$$
を $\theta$ に関して微分することに等価である。

Sinkhorn ダイバージェンスを含む汎関数の勾配を計算するコストは、汎関数自体の計算と同じである。このアプローチのいくつかの応用については~\citep{2016-bonneel-barycoord,2017-Genevay-AutoDiff}を参照されたい。また、コスト行列に関する Sinkhorn 反復の微分（Lagrange 的フレームワークで行われるもの）に関する初期の研究として、ランキング学習への応用を持つ~\citep{adams2011ranking}も参照されたい。
自動微分のさらなる詳細は~\citep{griewank2008evaluating,rall1981automatic,Neidinger10}に見出すことができ、特に勾配を計算する最速の方法である「逆モード」について述べられている。
実装の面では、すべての最近のディープラーニング Python フレームワークが最先端の逆モード微分と GPU/TPU 計算のサポートを備えており~\citep{theano2016,abadi2016tensorflow,pytorch}、Sinkhorn 損失の大規模応用には採用されるべきである。
このような自動微分技法の使用を強く推奨する。なぜなら、それらは~\eqref{eq-diff-loss-eul}および~\eqref{eq-diff-loss-lagr}の計算と同じ計算量を持ち、これらの公式は主に自動微分が何を計算しているかの理論的理解を得るために有用だからである。唯一の欠点は、逆モード自動微分がメモリを大量に消費することである（メモリは反復回数に比例して増加する）。しかし、この問題を軽減するサブサンプリング戦略が存在する~\citep{griewank1992achieving}。



%%%%%%%%%%%%%%%%%%%%%%%%%%%%%%%%%%%%%
\section{Wasserstein 重心、クラスタリング、および辞書学習}
\label{sec-bary}

教師なし学習における基本的な問題は、いくつかのデータ点の「平均」または「重心」を計算することである。距離空間 $(\X,d)$（ここで $d$ は距離またはより一般的にダイバージェンス）に属する点 $(x_s)_{s=1}^S \in \X^S$ のそのような加重平均を定義する古典的な方法は、変分問題
\eql{\label{eq-frechet-means}
	\umin{x \in \X} \sum_{s=1}^S \la_s d(x,x_s)^p
}
を与えられた重みの族 $(\la_s)_s \in \simplex_S$ に対して解くことである。ここで $p$ はしばしば $p=2$ に設定される。
$\X=\RR^d$ かつ $d(x,y)=\norm{x-y}_2$ の場合、これは $p=2$ では通常の線形平均 $x=\sum_s \la_s x_s$ の定義につながり、$p=1$ ではより高度な中央値点につながる。この形式を用いてさまざまな平均の概念（\eg $\X=\RR_+$ 上の調和平均や幾何平均）を得ることができる。
この過程はしばしば「Fr\'echet」平均または「Karcher」平均と呼ばれる（歴史的な説明は~\citet{karcher2014riemannian}を参照）。一般的な距離 $d$ に対して、問題~\eqref{eq-frechet-means}は通常困難な非凸最適化問題である。幸いなことに、optimal transport 距離の場合、この問題は存在が証明でき効率的な数値スキームが存在する凸計画として定式化できる。


%%%%%%%%%%%%
\paragraph{Wasserstein 空間上の Fr\'echet 平均.}

入力ヒストグラム $\{\b_s\}_{s=1}^S$（ここで $b_s \in \simplex_{n_s}$）と重み $\la \in \simplex_S$ が与えられたとき、Wasserstein 重心は以下を最小化することで計算される：
\eql{\label{eq-wass-discr}
	\umin{\a \in \simplex_n} \sum_{s=1}^S \la_s \MKD_{\C_s}(\a,\b_s),
}
ここでコスト行列 $\C_s \in \RR^{n \times n_s}$ を指定する必要がある。
典型的な設定は「Euler 的」であり、すべての重心が同じ格子上に定義され、$n_s=n$、$\C_s=\C=\distD^p$ は距離行列に設定され、
\eq{
	\umin{\a \in \simplex_n} \sum_{s=1}^S \la_s \WassD_p^p(\a,\b_s).
}
を解く。

重心問題~\eqref{eq-wass-discr}は、\citet{carlierekelandmatching}の先行するアイデアに基づき、\citet{Carlier_wasserstein_barycenter}によって任意の測度を含むより一般的な形で導入された。その提示は注意~\ref{rem-bary-carlier}に委ねる。
ヒストグラムに対する重心問題~\eqref{eq-wass-discr}は実際に線形計画問題である。なぜなら、各入力と重心自体との間の $S$ 個のカップリング $(\P_s)_s$ を探すことができ、構成上、それらは同じ行周辺を共有するように制約されるからである：
\eq{
	\umin{\a \in \simplex_n, (\P_s \in \RR^{n \times n_s})_s } \enscond{
		\sum_{s=1}^S \la_s \dotp{\P_s}{\C_s}
	}{
		\foralls s, \P_s^\top \ones_{n_s}=\a, \P_s^\top \ones_{n} = \b_s
	}.
}
この問題は線形計画問題であるが、そのスケールは中規模問題に対して汎用ソルバの使用を禁じる。双対に対する劣勾配降下法のような一次法に頼ることができる~\citep{Carlier-NumericsBarycenters}。

%%%%%%%%%%%%%%%%%%%%
\begin{rem2}{任意の測度の重心}\label{rem-bary-carlier}
	入力測度の集合 $(\be_s)_s$ がある空間 $\X$ 上に定義されている場合、重心問題は
	\eql{\label{eq-barycenter-generic}
		\umin{\al \in \Mm_+^1(\X)} \sum_{s=1}^S \la_s \MK_{\c}(\al,\be_s)
	}
	となる。
	$\X=\RR^d$ かつ $c(x,y)=\norm{x-y}^2$ の場合、\citet{Carlier_wasserstein_barycenter}は入力測度の1つが密度を持てば、この重心は一意であることを示した。
	問題~\eqref{eq-barycenter-generic}は、点 $(x_s)_{s=1}^S \in \X^S$ の重心を計算する問題を任意の測度に一般化したものと見なすことができる。実際、$\be_s=\de_{x_s}$ が単一の Dirac 質量であれば、~\eqref{eq-barycenter-generic}の解は $\de_{x^\star}$ であり、$x^\star$ は~\eqref{eq-frechet-means}を解く Fr\'echet 平均である。
	$c(x,y)=\norm{x-y}^2$ の場合、重心 $\al^\star$ の平均は必然的に平均の重心であることに注意する。\ie
	\eq{
		\int_\Xx x \d\al^\star(x) =  \sum_s \la_s \int_\Xx x \d\al_s(x),
	}
	また $\al^\star$ の台は $(\al_s)_s$ の台の凸包に含まれる。
	入力分布を離散分布で近似する際の無限次元最適化~\eqref{eq-barycenter-generic}の近似の整合性（すなわち代わりに~\eqref{eq-wass-discr}を解くこと）は~\citet{Carlier-NumericsBarycenters}で研究されている。
	また、~\eqref{eq-barycenter-generic}を多周辺 OT 問題として再定式化することも可能であることを注記する。注意~\ref{eq-multimarg-bary}を参照されたい。
\end{rem2}
%%%%%%%%%%%%%%%%%%%%

\begin{rem2}{Wasserstein 変分問題としての $k$-means}\label{rem-bary-kmeans}
入力測度の族 $(\be_s)_s$ がただ1つの測度 $\be$ に限定され、この測度が $\X=\RR^{\dim}$ の離散有限部分集合上に台を持ち、コストが二乗ユークリッド距離である場合、重心問題
	\eql{\label{eq-barycenter-kmeans}
		\umin{\al \in \Mm_{k}^1(\X)} \MK_{\c}(\al,\be),
	}
は、$\al$ が大きさ最大 $k$ の有限台を持つ離散測度に制約されている場合、$\be$ を取る通常の $k$-means 問題と等価であることを示すことができる。実際、$k$-means 問題が出力するセントロイドは解 $\al$ の台に対応し、その重みは各セントロイドに割り当てられた $\be$ 内の点の割合に対応することを容易に示せる。$\MK_{\c}$ をエントロピー正則化で近似すると、$k$-means のソフトクラスタリング変種（ガウス混合モデルなど~\citep{dessein2017parameter}）に現れる平滑化された割り当てが得られることを示すことができる。
\end{rem2}

%%%%%%%%%%%%%%%%%%%%
\begin{rem2}{分布の分布と整合性}
	~\eqref{eq-barycenter-generic}を、可能であれば無限の測度の集合に一般化することが可能である。この問題は、確率分布の空間 $\Mm_+^1(\X)$ 上の確率分布 $M$ を考えることで記述される。\ie $M \in \Mm_+^1(\Mm_+^1(\X))$。重心は
	\eql{\label{eq-bary-infinite}
		\umin{\al \in \Mm_+^1(\X)} \EE_M( \MK_{\c}(\al,\be) ) = \int_{\Mm_+^1(\X)} \MK_{\c}(\al,\be) \d M(\be),
	}
	の解であり、$\be$ は $M$ に従って分布するランダム測度である。$M$ に従って有限個 $S$ の入力測度 $(\be_s)_{s=1}^S$ を一様にランダムに抽出すれば、一様な重み $\la_s=1/S$ に対する~\eqref{eq-barycenter-generic}の解として $\hat \be_S$ を定義できる（ここで $\hat \be_S$ 自体がランダム測度であることに注意）。
	問題~\eqref{eq-barycenter-generic}は「離散的」な測度 $M=\sum_s \la_s \de_{\be_s}$ の特殊な場合に対応する。
	$\MK_{\c}(\hat\be_S,\al)$ の零への収束（期待値またはおおよその確率で。ここで $\al$ は~\eqref{eq-bary-infinite}の一意解）は重心の整合性に対応し、~\citep{BigotBarycenter,leGouic2016existence,bigot2012characterization}で証明されている。
	これは Wasserstein 空間上の大数の法則として解釈できる。この結果を中心極限定理に拡張することは重要な問題である。この問題の最近の定式化と特定の場合（1次元分布とガウス測度）における解については~\citep{panaretos2016amplitude}および~\citep{agueh2017vers}を参照されたい。
\end{rem2}
%%%%%%%%%%%%%%%%%%%%


%%%%%%%%%%%%%%%%%%%%
\begin{rem2}{不動点写像}
ユークリッド空間 $\X=\RR^\dim$ で ground cost $\c(x,y)=\norm{x-y}^2$ を扱う場合、輸送写像を用いて重心問題を研究することが可能である。実際、$\al$ が密度を持てば、注意~\ref{rem-exist-mongemap}により、$\al$ と $\al_s$ の間の最適輸送写像 $\T_s$、特に $\T_{s,\sharp}\al = \al_s$ を満たすものを定義できる。平均写像
\eq{
	\T^{(\al)} \eqdef \sum_{s=1}^S \la_s \T_s
}
（上記の記法は $\al$ へのこの写像の依存性を明示している）は、それ自体が $\al$ と $\T^{(\al)}_\sharp \al$ の間の最適写像である（最適写像の正の組み合わせは、Brenier の定理（注意~\ref{rem-exist-mongemap}）により凸関数の勾配の和に等しく、凸関数の和の勾配に等しいので、Brenier の定理により再び最適である）。
\citet{Carlier_wasserstein_barycenter}が示したように、重心問題~\eqref{eq-bary-infinite}の一次最適性条件は実際に最適測度 $\al^\star$（重心）において $\T^{(\al^\star)}=\Identity_{\RR^\dim}$（恒等写像）と読め、~\citep{alvarez2016fixed}では重心 $\al^\star$ が不動点方程式
\eql{\label{eq-fixed-point-bary}
	G(\al) = \al
	\qwhereq
	G(\al) \eqdef \T^{(\al)}_\sharp \al,
}
の一意解（~\citep[Theo. 2]{zemel2017fr}で明確にされた正則性条件の下で）であることが示されている。
入力測度に対する穏やかな条件の下で、\citet{alvarez2016fixed}および\citet{zemel2017fr}は、$\al$ が重心でなければ $\al \mapsto G(\al)$ が~\eqref{eq-bary-infinite}の目的関数を狭義に減少させること、および不動点反復 $\itt{\al} \eqdef G( \it{\al} )$ が重心 $\al^\star$ に収束することを示した。
この不動点アルゴリズムは、最適輸送写像が閉じた形で既知の場合（\eg ガウス分布の場合）に用いることができる。このアルゴリズムを同じサイズの経験測度に適用すると、Monge 写像の代わりに最適割り当ての計算となる。より一般的な任意のサイズの離散測度に対しても、重心射影~\eqref{eq-baryproj}を用いてスキームを適応させることができる~\citep{CuturiBarycenter}。
\end{rem2}
%%%%%%%%%%%%%%%%%%%%



%%%%%%%%%%%%
\paragraph{特殊な場合.}

一般に、~\eqref{eq-wass-discr}や~\eqref{eq-barycenter-generic}を解くことは単純でないが、解が明示的または単純であるいくつかの特殊な場合が存在する。


%%%%%%%%%%%%%%%%%%%%
\begin{rem1}{ガウス分布の重心}
	\citep{Carlier_wasserstein_barycenter}において、二乗ユークリッドコスト $c(x,y)=\norm{x-y}^2$ に対するガウス分布 $\al_s = \Nn(\mean_s,\cov_s)$ の重心は、それ自体がガウス分布 $\Nn(\mean^\star,\cov^\star)$ であることが示されている。
	~\eqref{eq-dist-gauss}を用いると、重心の平均は入力の平均であり
	\eq{
		\mean^\star = \sum_s \la_s \mean_s
	}
	共分散は
	\eq{
		\umin{ \cov } \sum_s \la_s \Bb(\cov,\cov_s)^2,
	}
	を最小化する。ここで $\Bb$ は Bures 計量~\eqref{eq-bure-defn}である。\citep{Carlier_wasserstein_barycenter}で研究されたように、この凸問題の一次最適性条件は、$\cov^\star$ が写像の一意の正定値不動点であることを示す：
	\eq{
		\cov^\star = \Psi(\cov^\star)
		\qwhereq
		\Psi(\cov) \eqdef \sum_s \la_s ( \cov^{\frac{1}{2}} \cov_s \cov^{\frac{1}{2}}  )^{\frac{1}{2}},
	}
	ここで $\cov^{\frac{1}{2}}$ は半正定値行列の平方根である。
	この結果は~\citep{KnottSmith,RuschendorfUckelmann}から知られており、~\citep{Carlier_wasserstein_barycenter}で証明されている。
	$\Psi$ は狭義縮小ではないが、この不動点写像を反復すること、\ie $\itt{\cov} \eqdef \Psi( \it{\cov} )$ を定義することは、実際には解 $\cov^\star$ に収束する。
	この方法はテクスチャ合成~\citep{2014-xia-siims}に応用されている。\citet{alvarez2016fixed}はまた、別の写像
	\eq{
		\bar\Psi(\cov) \eqdef
		\cov^{-\frac{1}{2}}
		\Big( \sum_s \la_s ( \cov^{\frac{1}{2}} \cov_s \cov^{\frac{1}{2}}  )^{\frac{1}{2}} \Big)^2
		\cov^{-\frac{1}{2}}
	}
	の使用を提案しており、反復 $\itt{\cov} \eqdef \bar\Psi( \it{\cov} )$ は収束する。
	これは~\eqref{eq-fixed-point-bary}で定義された不動点写像 $G$ がガウス分布を保存するためであり、実際に
	\eq{
		G( \Nn(\mean,\cov) ) = \Nn(\mean^\star,\bar\Psi(\cov)).
	}
	が成り立つ。
	図~\ref{fig-bary-gaussian}は、4つの2次元ガウス分布間の重心計算の2つの例を示している。
\end{rem1}
%%%%%%%%%%%%%%%%%%%%



\begin{figure}[h!]
\centering
\begin{tabular}{@{}c@{\hspace{20mm}}c@{}}
\includegraphics[width=.3\linewidth]{barycenters-gaussians/tensor-interp-70} &
\includegraphics[width=.3\linewidth]{barycenters-gaussians/tensor-interp-98}
\end{tabular}
\caption{\label{fig-bary-gaussian}
2次元における4つのガウス分布間の重心。各ガウス分布は共分散の主軸に沿った楕円で表示され、対応する固有値に比例した伸びを持つ。
}
\end{figure}


%%%%%%%%%%%%%%%%%%%%
\begin{rem1}{1次元の場合}\label{rem-bary-1d}
	1次元分布の場合、$\Wass_p$ 重心は、注意~\ref{rem-1d-ot-generic}で詳述されるように、輸送が単調再配置であるという事実を用いてほぼ閉じた形で計算できる。
	最も単純な場合は $n$ 個の点を持つ経験測度の場合であり、\ie $\be_s = \frac{1}{n}\sum_{i=1}^n \de_{y_{s,i}}$、ここで点はソートされていると仮定する：$y_{s,1} \leq y_{s,2} \leq \ldots$。~\eqref{eq-1d-empirical}を用いると、重心 $\al_\la$ も $n$ 個の点上の経験測度であり
	\eq{
		\al_\la = \frac{1}{n}\sum_{i=1}^n \de_{x_{\la,i}}
		\qwhereq
		x_{\la,i} = A_\la(x_{s,i})_s,
	}
	ここで $A_\la$ は重心写像
	\eq{
	 	A_\la(x_s)_s \eqdef \uargmin{x \in \RR} \sum_{s=1}^S \la_s |x-x_{s}|^p
	}
	である。
	例えば、$p=2$ の場合、$x_{\la,i} = \sum_{s=1}^S \la_s x_{s,i}$ となる。
	一般の場合、~\eqref{eq-cumul-defn}で定義される累積分布関数を用い、~\eqref{eq-wass-cumul}を用いて
	\eq{
		\foralls r \in [0,1], \quad
		\cumul{\al_\la}^{-1}(r) = A_\la( \cumul{\al_s}^{-1}(r) )_{s=1}^S,
	}
	が成り立つ。これは例えば、単純なソート手続きを用いて $n$ 個未満の点上に台を持つ離散測度間の重心を $O(n\log(n))$ 演算で計算するために使用できる。
\end{rem1}
%%%%%%%%%%%%%%%%%%%%



%%%%%%%%%%%%%%%%%%%%
\begin{rem1}{単純な場合}
	$\T_{r,u} : x \mapsto rx+u$ をスケーリングと平行移動とし、$\al_{s} = T_{r_s,u_s,\sharp} \al_0$ が初期テンプレート測度のスケーリングと平行移動により得られると仮定すると、重心 $\al_\la$ もスケーリングと平行移動により得られる：
	\eq{
		\al_\la = T_{r^\star,u^\star,\sharp} \al_0
		\qwhereq
		\choice{
			r^\star = (\sum_s \la_s/r_s)^{-1},  \\
			u^\star = \sum_s \la_s u_s.
		}
	}
\end{rem1}
%%%%%%%%%%%%%%%%%%%%


%%%%%%%%%%%%%%%%%%%%
\begin{rem1}{$S=2$ の場合}
	$\X=\RR^d$ かつ $c(x,y)=\norm{x-y}^2$ の場合（これはより一般的に測地的空間に拡張できる）、$S=2$ 個の測度 $(\al_0,\al_1)$ 間の重心は、~\eqref{eq-mc-cann-interp}で既に導入された McCann 補間である。$\T_\sharp \al_0=\al_1$ を Monge 写像とすると、重心 $\al_\la$ は $\al_\la = (\la_1 \Id + \la_2 \T)_{\sharp} \al_0$ と読める。
	公式~\eqref{eq-mccann-discrete}は離散の場合の計算方法を説明する。
\end{rem1}
%%%%%%%%%%%%%%%%%%%%



%%%%%%%%%%%%
\paragraph{重心のエントロピー近似.}

エントロピー平滑化を用い、ある $\epsilon>0$ に対して~\eqref{eq-wass-discr}の解を以下で近似できる：
\eql{\label{eq-entropic-bary}
	\umin{\a \in \simplex_n} \sum_{s=1}^S \la_s \MKD_{\C_s}^\epsilon(\a,\b_s)
}
これは滑らかな凸最小化問題であり、勾配降下法~\citep{CuturiBarycenter,GramfortPC15}で解くことができる。代替として、半双対に対する降下法（典型的には準 Newton 法）~\citep{2016-Cuturi-siims}を用いることができ、これは重心に追加の正則化を課す（例えば与えられたノルムに関する滑らかさを課す）際に有用である。
\citet{2015-benamou-cisc}が注目したように、より単純で非常に効果的なアプローチは~\eqref{eq-entropic-bary}を（重み付き）KL 射影問題として書き直すことである：
\eql{\label{eq-bary-entropy-couplings}
	\umin{ (\P_s)_s } \enscond{ \sum_{s} \la_s \epsilon \KLD( \P_s|\K_s ) }{
		\foralls s, \transp{\P_s}\ones_m = \b_s, \:
		\P_1\ones_1 =  \cdots = \P_S\ones_S,
	}
}
ここで $\K_s \eqdef e^{-\C_s/\epsilon}$ と記した。ここで、重心 $\a$ はすべてのカップリング $\P_s \in \RR^{n \times n_s}$ の行周辺として $\a = \P_1\ones_1 =  \cdots = \P_S\ones_S$ と暗黙的にコード化されている。
\citet{2015-benamou-cisc}が詳述するように、この問題に Sinkhorn を一般化でき、これは反復射影にも対応する。これは第~\ref{sec-generalized}節で詳述される一般化 Sinkhorn の特殊な場合と見なすこともできる。
~\eqref{eq-bary-entropy-couplings}を解く最適カップリング $(\P_s)_s$ はスケーリング形式で
\eql{\label{eq-bary-opt}
	\P_s=\diag(\uD_s)\K\diag(\vD_s),
}
と計算され、スケーリングは逐次的に
\begin{align}\label{eq-sinkhorn-bary}
	\foralls s \in \range{1,S}, \quad \itt{\vD}_s &\eqdef \frac{\b_s}{\transp{\K}_s \it{\uD}_s}, \\
	\foralls s \in \range{1,S}, \quad  \itt{\uD}_s &\eqdef \frac{\itt{\a}}{\K_s \itt{\vD}_s}, \label{eq-sinkhorn-bary-2}\\
		\qwhereq
		\itt{\a} &\eqdef \prod_s (  \K_s \itt{\vD}_s )^{\la_s}. \label{eq-sinkhorn-bary-3}
\end{align}
と更新される。これらの反復を導出する別の方法は、双対問題の変数に対して交互最小化を行うことであり、以下の命題で詳述される。

\begin{prop}
	~\eqref{eq-bary-opt}に現れる最適な $(\uD_s,\vD_s)$ は $(\uD_s,\vD_s) = ( e^{\fD_s/\epsilon},e^{\gD_s/\epsilon} )$ と書くことができる。ここで $( \fD_s,\gD_s )_s$ は以下の計画（その値は~\eqref{eq-entropic-bary}の値と一致する）の解である：
	\eql{\label{eq-dual-bary-entropy}
		\umax{ ( \fD_s,\gD_s )_s } \enscond{
		\sum_s \la_s \pa{
			\dotp{\gD_s}{\b_s} - \epsilon \dotp{\K_s e^{\gD_s/\epsilon}}{e^{\fD_s/\epsilon}}
		}
		}{
			\sum_s \la_s \fD_s = 0
		}.
	}
\end{prop}

\begin{proof} ~\eqref{eq-bary-entropy-couplings}に Lagrange 乗数を導入すると
	\begin{align*}
		\umin{ (\P_s)_s, \a }
		\umax{ ( \fD_s,\gD_s )_s }
		\sum_{s} \la_s \Big(
			\epsilon \KLD( \P_s|\K_s )
			+
			\dotp{\a - \P_s\ones_m}{ \fD_s } \\
		\qquad\qquad	+
			\dotp{\b_s - \transp{\P_s}\ones_m}{ \gD_s }
		\Big).
	\end{align*}
	が得られる。強双対性が成り立つので、min と max を交換でき、
		\begin{align*}
		&\umax{ ( \fD_s,\gD_s )_s }
			\sum_s \la_s \pa{
				 \dotp{\gD_s}{\b_s}
		+
		\umin{\P_s} \epsilon \KLD( \P_s|\K_s ) - \dotp{\P_s}{\fD_s \oplus \gD_s }
		} \\
	&\qquad\qquad	+ \umin{\a}
			\dotp{ \sum_{s} \la_s \fD_s }{\a}.
	\end{align*}
	が得られる。$\a$ に関する明示的最小化は、制約 $\sum_{s} \la_s \fD_s=0$ を与え、
	\begin{align*}
		\umax{ ( \fD_s,\gD_s )_s }
			\sum_s \la_s
				 \dotp{\gD_s}{\b_s}
		-
		\epsilon \KLD^*\pa{ \frac{\fD_s \oplus \gD_s}{\epsilon}|\K_s },
	\end{align*}
	が得られる。ここで $\KLD^*(\cdot|\K_s)$ は関数 $\KLD^*(\cdot|\K_s)$ の Legendre 変換~\eqref{eq-legendre}である。
	この Legendre 変換は
	\eql{\label{eq-legendre-kl}
		\KLD^*(\VectMode{U}|\K) = \sum_{i,j} \K_{i,j} (e^{\VectMode{U}_{i,j}}-1),
	}
	となり、所望の公式を示す。
	~\eqref{eq-legendre-kl}を示すために、この関数は分離可能であるから、
	\eq{
		\foralls (u,k) \in \RR_+^2, \quad
		\KLD^*(u|k) \eqdef \umax{r} ur -   \pa{ r \log(r/k) - r + k  }
	}
	を計算する必要があり、その最適性条件は $u=\log(r/k)$、\ie $r = k e^{u}$ であるから、結果が得られる。
\end{proof}

各 $\gD_s$ に関して他のすべての変数を固定したまま~\eqref{eq-dual-bary-entropy}を最小化することは~\eqref{eq-sinkhorn-bary}により閉じた形で得られる。
すべての $(\fD_s)_s$ に関して~\eqref{eq-dual-bary-entropy}を最小化するには~\eqref{eq-sinkhorn-bary-3}を用いて $\a$ を解く必要があり、式~\eqref{eq-sinkhorn-bary-2}の表現に至る。



図~\ref{fig-barycenters-images}および図~\ref{fig-barycenters-shapes}は2次元および3次元の形状補間への応用を示す。図~\ref{fig-barycenters-surfaces}は曲面上での重心の計算を示しており、ground cost は測地距離の二乗である。この図では、注意~\ref{rem-geod-heat}で詳述される heat 法による測地距離近似を用いて計算が行われている。詳細およびコンピュータグラフィックスおよびイメージングサイエンスへの他の応用については~\citep{2015-solomon-siggraph}を参照されたい。

\newcommand{\FigBaryImA}[2]{\includegraphics[width=.088\linewidth]{barycenters-2d/annulus-cross-heart-2disk/shape-#1-#2}}
\newcommand{\FigBaryImLineA}[1]{\FigBaryImA{#1}{1} & \FigBaryImA{#1}{2} & \FigBaryImA{#1}{3} & \FigBaryImA{#1}{4} & \FigBaryImA{#1}{5}}
\newcommand{\FigBaryImB}[2]{\includegraphics[width=.088\linewidth]{barycenters-2d/cat-star8-thinspiral-trefle/shape-#1-#2}}
\newcommand{\FigBaryImLineB}[1]{\FigBaryImB{#1}{1} & \FigBaryImB{#1}{2} & \FigBaryImB{#1}{3} & \FigBaryImB{#1}{4} & \FigBaryImB{#1}{5}}


\begin{figure}[h!]
\centering
\fbox{
\begin{tabular}{@{}c@{}c@{}c@{}c@{}c@{}}
\FigBaryImLineA{1} \\
\FigBaryImLineA{2} \\
\FigBaryImLineA{3} \\
\FigBaryImLineA{4} \\
\FigBaryImLineA{5}
\end{tabular}
}
\fbox{
\begin{tabular}{@{}c@{}c@{}c@{}c@{}c@{}}
\FigBaryImLineB{1} \\
\FigBaryImLineB{2} \\
\FigBaryImLineB{3} \\
\FigBaryImLineB{4} \\
\FigBaryImLineB{5}
\end{tabular}
}
\caption{\label{fig-barycenters-images}
エントロピー正則化~\eqref{eq-entropic-bary}を用いた4つの入力2次元形状間の重心。二値形状を表示するために、表示画像は閾値処理された密度を示している。
重み $(\la_s)_s$ は正方形の4つの角に関して双線形である。
}
\end{figure}



\newcommand{\FigBaryShapeA}[2]{\includegraphics[width=.09\linewidth,trim=130 78 110 60,clip]{barycenters-shapes/duck-spiky-moomoo_s0-double-torus/barycenter-#1-#2}}
\newcommand{\FigBaryShapeLineA}[1]{\FigBaryShapeA{#1}{0} & \FigBaryShapeA{#1}{1} & \FigBaryShapeA{#1}{2} & \FigBaryShapeA{#1}{3} & \FigBaryShapeA{#1}{4}}
\newcommand{\FigBaryShapeB}[2]{\includegraphics[width=.087\linewidth,trim=160 65 150 55,clip]{barycenters-shapes/mushroom-torus-hand1-trim-star/barycenter-#1-#2}}
\newcommand{\FigBaryShapeLineB}[1]{\FigBaryShapeB{#1}{0} & \FigBaryShapeB{#1}{1} & \FigBaryShapeB{#1}{2} & \FigBaryShapeB{#1}{3} & \FigBaryShapeB{#1}{4}}

\begin{figure}[h!]
\centering
\fbox{
\begin{tabular}{@{}c@{}c@{}c@{}c@{}c@{}}
\FigBaryShapeLineA{0} \\
\FigBaryShapeLineA{1} \\
\FigBaryShapeLineA{2} \\
\FigBaryShapeLineA{3} \\
\FigBaryShapeLineA{4}
\end{tabular}
}
\fbox{
\begin{tabular}{@{}c@{}c@{}c@{}c@{}c@{}}
\FigBaryShapeLineB{0} \\
\FigBaryShapeLineB{1} \\
\FigBaryShapeLineB{2} \\
\FigBaryShapeLineB{3} \\
\FigBaryShapeLineB{4}
\end{tabular}
}
\caption{\label{fig-barycenters-shapes}
エントロピー正則化~\eqref{eq-entropic-bary}を用いた4つの入力3次元形状間の重心。重み $(\la_s)_s$ は正方形の4つの角に関して双線形である。
形状は、形状の境界内で一様であり外側では零である測度として表現される。
}
\end{figure}
Wasserstein 重心の効率的な計算は現時点でなお活発な研究テーマである~\citep{NIPS2017_6858,NIPS2018_8274}。方法論的関心に加えて、Wasserstein 重心は形状解析の分野を超えて多くの応用を見出している。画像処理~\citep{rabin-ssvm-11}、特に色修正~\citep{2015-solomon-siggraph}（図~\ref{fig-colors}を参照）、測度を要約するための Bayes 計算~\citep{srivastava2015scalable,srivastava2015wasp}、入力測度を他の既知の測度の Wasserstein 重心として表現する非線形次元削減~\citep{2016-bonneel-barycoord}に用いられている。これらの問題はすべて、自動微分（注意~\ref{rem-auto-diff}を参照）を用いて正確に最適化できる複雑な非凸目的関数をもたらす。
重心の計算と密接に関連する問題には、Wasserstein 空間上の主成分分析の計算（例えば~\citep{SeguyCuturi,bigot2017geodesic}を参照）およびテンプレートモデルの統計的推定~\citep{boissard2015distribution}がある。重心を計算する能力は、確率測度の空間上の $k$-means~\citep{del2016robust,ho2017multilevel}のようなより高度なクラスタリング手法を可能にする。




\begin{figure}[h!]
\centering
\begin{tabular}{@{}c@{}c@{}c@{}c@{}c@{}}
\includegraphics[width=.195\linewidth]{barycenters-surfaces/scisor-1}&
\includegraphics[width=.195\linewidth]{barycenters-surfaces/scisor-2}&
\includegraphics[width=.195\linewidth]{barycenters-surfaces/scisor-3}&
\includegraphics[width=.195\linewidth]{barycenters-surfaces/scisor-4}&
\includegraphics[width=.195\linewidth]{barycenters-surfaces/scisor-5}
\end{tabular}
\caption{\label{fig-barycenters-surfaces}
heat 法による高速カーネル近似（注意~\ref{rem-geod-heat}を参照）を用いて計算された、曲面上の2つの入力測度間の重心補間。~\citep{2015-solomon-siggraph}より抜粋。
}
\end{figure}



\begin{figure}[h!]
\centering
\includegraphics[width=1\linewidth]{colors/colors-bary}
\caption{\label{fig-colors}
2つの入力画像の3次元色経験ヒストグラム間の補間（簡単のため、ここでは2次元の彩度射影のみを可視化）。修正されたヒストグラムは、注意~\ref{rem-barycenric-proj}で詳述される重心射影を用いて入力画像に適用される。~\citep{2015-solomon-siggraph}より抜粋。
}
\end{figure}



%%%%%%%%%%%%%%%%%%%%
\begin{rem1}{Wasserstein 伝播}
	\citet{Solomon-ICML}で研究されたように、重心問題~\eqref{eq-wass-discr}を一般化することが可能であり、グラフ $\Gg$ における与えられたノード集合 $U$ での分布 $(\b_u)_{u \in U}$ を、ノードの補集合 $V$ 上の固定された入力分布 $(\b_v)_{v \in V}$ が与えられたもとで探索する。未知量は、グラフの辺を形成するすべてのノード対 $(r,s) \in \Gg$ 間の総輸送距離を最小化することで決定される：
	\eql{\label{eq-w-propag}
		\umin{(\b_u  \in \simplex_{n_u})_{u \in U}}
			\sum_{(r,s) \in \Gg}  \MKD_{\C_{r,s}}(\b_r,\b_s),
	}
	ここでコスト行列 $\C_{r,s} \in \RR^{n_r \times n_s}$ はユーザが指定する必要がある。
	重心問題~\eqref{eq-wass-discr}はこの問題の特殊な場合であり、考えるグラフ $\Gg$ が「星型」で、$U$ が他のすべての頂点 $V$ に接続された単一の頂点である場合に対応する（$\b_s$ に付随する重み $\la_s$ はコスト行列に吸収できる）。
	各辺 $(r,s) \in \Gg$ に対してカップリング $\P_{r,s} \in \CouplingsD( \b_r,\b_s )$ を明示的に導入し、エントロピー正則化を用いると、~\eqref{eq-bary-entropy-couplings}と同様にこの問題を書き直すことができ、Sinkhorn 反復~\eqref{eq-sinkhorn-bary}をこの問題に拡張できる（これは第~\ref{sec-generalized}節で詳述される一般化 Sinkhorn アルゴリズムの形に問題を再定式化することでも導出できる）。
	グラフ上のこの離散変分問題~\eqref{eq-w-propag}は、グラフを連続領域で置き換えたときの Dirichlet エネルギーの定義に一般化できる~\citep{solomon2013dirichlet}。これにより、画像および曲面処理に応用を見出す測度値調和関数の定義につながる。理論的分析については~\citet{Lavenant2017}を、非二次（全変動）汎関数への拡張およびイメージングへの応用については~\citet{vogt2017measure}も参照されたい。
\end{rem1}
%%%%%%%%%%%%%%%%%%%%



%%%%%%%%%%%%%%%%%%%%%%%%%%%%%%%%%%%%%%%%%%%%%%%%%%%%%%%%%%%%%%%%%%%%%%%%%%
%%%%%%%%%%%%%%%%%%%%%%%%%%%%%%%%%%%%%%%%%%%%%%%%%%%%%%%%%%%%%%%%%%%%%%%%%%
%%%%%%%%%%%%%%%%%%%%%%%%%%%%%%%%%%%%%%%%%%%%%%%%%%%%%%%%%%%%%%%%%%%%%%%%%%
\section{勾配流}
\label{sec-grad-flows}

滑らかな関数 $\a \mapsto F(\a)$ が与えられたとき、標準的な勾配降下法
\eql{\label{eq-explicit-euclidean}
	\itt{\a} \eqdef \it{\a} - \tau \nabla F(\it{\a}),
}
を用いることができる。ここで $\tau$ は十分小さなステップサイズである。これはいわゆる「陽的」最小化スキームに対応し、滑らかな関数 $F$ に対してのみ適用される。非滑らかな関数に対しては、代わりに「陰的」スキームを用いることができ、これは近接点アルゴリズムとも呼ばれる（例えば~\citet{BauschkeCombettes11}を参照）：
\eql{\label{eq-implicit-euclidean}
	\itt{\a} \eqdef \Prox_{\tau F}^{\norm{\cdot}}(\it{\a}) \eqdef \uargmin{\a} \frac{1}{2}\norm{\a-\it{\a}}^2 + \tau F(\a).
}
これは~\eqref{eq-prox-eucl}で既に遭遇したユークリッド近接作用素に対応することに注意する。
更新~\eqref{eq-explicit-euclidean}は陽的作用素 $\Id-\tau\nabla F$ の反復と理解でき、~\eqref{eq-implicit-euclidean}は陰的作用素 $(\Id+\tau\nabla F)^{-1}$ を使用する。凸な $F$ に対して、反復~\eqref{eq-implicit-euclidean}は $\tau>0$ の任意の値に対して常に収束する。

関数 $F$ がヒストグラムの単体 $\simplex_n$ 上に定義されている場合、~\eqref{eq-implicit-euclidean}における $\ell^2$ ノルム $\norm{\cdot}$ の代わりに optimal transport 距離を用いることは理にかなっており、
\eql{\label{eq-grad-flow-discr}
	\itt{\a} \eqdef \uargmin{\a} \WassD_p(\a,\it{\a})^p + \tau F(\a)
}
を解く。

%%%%%%%%%%%%%%%%%%
\begin{rem2}{Wasserstein 勾配流}
式~\eqref{eq-grad-flow-discr}は反復
\eql{\label{eq-grad-flow-stepping-cont}
	\itt{\al} \eqdef \uargmin{\al} \Wass_p(\al,\it{\al})^p + \tau F(\al)
}
を定義することで任意の測度に一般化できる。ここで $F$ は $\Mm_+^1(\X)$ 上に定義された関数である。この陰的時間ステッピングは連続的な流れを構築するための有用なツールであり、形式的に極限 $\tau \rightarrow 0$ を取り時間 $t=\tau \ell$ を導入すると、$\it{\al}$ は連続流 $t \in \RR_+ \mapsto \al_t$ を近似することが意図される。特殊な場合 $p=2$ かつ $\X=\RR^d$ では、形式的な計算により $\al_t$ は
\eql{\label{eq-grad-flow-cont}
	\pd{\al_t}{t} = \diverg( \al_t \nabla( F'(\al_t) ) ),
}
の形の PDE を解くことが期待される。ここで $F'(\al)$ は、連続関数 $F'(\al) \in \Cc(\X)$ であって
\eq{
	F(\al+\epsilon\xi) = F(\al) + \epsilon \int_\X F'(\al) \d\xi(x) + o(\epsilon)
}
を満たすという意味での関数 $F$ の導関数を表す。
典型的な例は $F=-\H$ を用いる場合であり、$\H(\al)=\KL(\al|\Ll_{\RR^d})$ は $\X=\RR^\dim$ 上の Lebesgue 測度 $\Ll_{\RR^d}$ に関する相対エントロピーである：
\eql{\label{eq-entropy-cont}
	\H(\al) = - \int_{\RR^\dim} \density{\al}(x) (\log(\density{\al}(x)) - 1) \d x
}
（$\al$ が密度を持たないとき $\H(\al)=-\infty$ と設定する）。この場合~\eqref{eq-grad-flow-cont}は、この負エントロピーの勾配流が線形熱拡散
\eql{\label{eq-heat}
	\pd{\al_t}{t} = \Delta \al_t,
}
であることを示す。ここで $\Delta$ は空間 Laplacian である。
したがって、熱拡散は Dirichlet エネルギー $\int_{\RR^\dim} \|\nabla \density{\al}(x)\|^2 \d x$ の「古典的な」ユークリッド流（質量振幅に関するある種の「垂直」運動）と解釈することもでき、あるいは optimal transport 流に対するエントロピー（質量位置に関するある種の「水平」運動）とも解釈できる。
Wasserstein 勾配流への関心は Jordan, Kinderlehrer, Otto の先駆的論文~\citep{jordan1998variational}により火がつけられ、これらの発展は彼らの研究にちなんで「JKO 流」と呼ばれることが多い。\citet{ambrosio2006gradient}のモノグラフで詳細に示されるように、JKO 流は距離空間における勾配流の特殊な場合である。最近のサーベイ論文~\citep{santambrogio2017euclidean}も参照されたい。
JKO 流は、特に多孔質媒体方程式~\citep{otto2001geometry}、全変動流~\citep{carlier2017total}、量子ドリフト~\citep{GianazzaARMA}、多様体上の熱発展~\citep{ErbarHeatManifold}などの非線形発展方程式を研究するために用いることができる。その柔軟な形式は、\citeauthor{maury2010macroscopic}が群衆運動~\citep{maury2010macroscopic}をモデル化するために用いた混雑制約（任意の点での密度の上界）のような解への制約を許容する（レビュー論文~\citep{SantambrogioCrowdReview}も参照）。
\end{rem2}
%%%%%%%%%%%%%%%%%%%


%%%%%%%%%%%%%%%%%%
\begin{rem2}{距離空間における勾配流}
陰的ステッピング~\eqref{eq-grad-flow-stepping-cont}は、~\citep{ambrosio2006gradient}で詳述される距離空間 $(\Xx,\dist)$（ここで $\dist$ は距離）上の勾配流を定義するためのより一般的な形式の特殊な場合である。
$x \in \Xx$ に対して定義されたある関数 $F(x)$ に対して、陰的離散最小化ステップは
\eql{\label{eq-implicit-metricflow}
	\itt{x} \in \uargmin{x \in \Xx} \dist(\it{x},x)^2 + \tau F(x)
}
と定義される。
JKO ステップ~\eqref{eq-grad-flow-stepping-cont}は確率分布の空間上の Wasserstein 距離の使用に対応する。いくつかの場合、$\tau \rightarrow 0$ かつ $k\tau=t$ のとき~\eqref{eq-implicit-metricflow}が連続流の極限 $x_t$ を許容することを示すことができる。
$\Xx$ がユークリッド構造も持つ場合、$F$ を線形化することで陽的ステッピングが定義される：
\eql{\label{eq-explicit-metricflow}
	\itt{x} = \uargmin{x \in \Xx} \dist(\it{x},x)^2 + \tau \dotp{ \nabla F(\it{x}) }{ x }.
}
陰的公式~\eqref{eq-implicit-metricflow}とは対照的に、これは通常計算が容易であるが不安定になり得る。陰的ステップは常に安定であり、非滑らかな $F$ に対しても定義されるが、通常閉じた形では得られない。
図~\ref{fig-gradflow-metric}は、関数 $F(x) = \norm{x}^2$ を $\Xx=\RR^2$ 上で、距離 $d(x,y)=\norm{x-y}_p=(|x_1-y_1|^p+|x_2-y_2|^p)^{\frac{1}{p}}$ に対して $p$ のいくつかの値についてこの概念を示している。
陽的スキーム~\eqref{eq-explicit-metricflow}は $p=1$ と $p=+\infty$ で不安定であり、$p=1$ では軸に沿ったステップ（座標ごとの降下法）を与える。
対照的に、陰的スキーム~\eqref{eq-implicit-metricflow}は安定である。特に $p=1$ の場合、2つの座標が等しいとき、次のステップが対角方向に作用することに注意する。
\end{rem2}
%%%%%%%%%%%%%%%%%%


\begin{figure}[h!]
\centering
\begin{tabular}{cc}
\includegraphics[width=.45\linewidth]{gradflow-metric/explicit}&
\includegraphics[width=.45\linewidth]{gradflow-metric/implicit}\\
Explicit & Implicit
\end{tabular}
\caption{\label{fig-gradflow-metric}
関数 $f(x) = \norm{x}^2$ を $\Xx=\RR^2$ 上で最小化するための、距離 $d(x,y)=\norm{x-y}_p$ に対する $p$ のいくつかの値についての陽的および陰的勾配流の比較。
}
\end{figure}


%%
\begin{rem1}{粒子系を用いた Lagrange 的離散化}\label{rem-lagr-discretization}
~\eqref{eq-grad-flow-discr}の有限次元問題は、測度の空間上の流れ~\eqref{eq-grad-flow-stepping-cont}の Euler 的離散化と解釈できる。
問題を離散化する別の方法は、いわゆる粒子系を用いた Lagrange 的方法であり、解を時間とともに移動する（離散的な）経験測度としてパラメータ化し、その測度の位置（重みではなく）が関心のある変数となる。実際には、時間でインデックスされた動的な点群の粒子 $\al_t = \frac{1}{n}\sum_{i=1}^n \de_{x_i(t)}$ を考えることができる。
元の問題~\eqref{eq-grad-flow-discr}は、点 $X(t)=(x_i(t))_i \in \X^n$ の動力学を規定する $n$ 個の連立 ODE で置き換えられる。
エネルギー $F$ が離散測度に対して有限であれば、単純に $\Ff(X)=F( \frac{1}{n}\sum_{i=1}^n \de_{x_i} )$ と定義できる。典型的な例は線形関数 $F(\al) = \int_\X V(x) \d\al(x)$ と二次相互作用 $F(\al) = \int_{\X^2} W(x,y) \d\al(x)\d\al(y)$ であり、その場合それぞれ
\eq{
	\Ff(X) = \frac{1}{n} \sum_i V(x_i) \qandq \Ff(X) = \frac{1}{n^2} \sum_{i,j} W(x_i,x_j)
}
を用いることができる。
密度を持つ測度に対してのみ有限である一般化エントロピーのような関数に対しては、点群を密度に変換するための密度推定器を適用する必要があり、これにより $n \rightarrow +\infty$ のとき $F$ と整合的な関数 $\Ff(x)$ を定義することもできる。典型的な例は~\eqref{eq-entropy-cont}で定義されるエントロピー $F(\mu) = \H(\al)$ であり、（定数項を除いて）整合的な推定器は最近傍距離の対数の総和により得られる：
\eql{\label{eq-entropy-discretized}
	\Ff(X) = \frac{1}{n}\sum_{i} \log( d_X(x_i) )
	\qwhereq d_X(x) = \umin{x' \in X, x' \neq x} \norm{x-x'};
}
ノンパラメトリックエントロピー推定器のレビューについては~\citet{beirlant1997nonparametric}を参照されたい。
十分小さなステップサイズ $\tau$ に対して、$\X=\RR^\dim$ と仮定すると、Wasserstein 距離 $\Wass_2$ は点上のユークリッド距離と一致する。\ie $|t-t'|$ が十分小さければ $\Ww_2(\al_t,\al_{t'})=\norm{X(t)-X(t')}$。したがって勾配流は位置上のユークリッド流 $X'(t) = -\nabla \Ff(X(t))$ と等価であり、時刻 $t_k=\tau k$ において~\eqref{eq-explicit-euclidean}と同様に陽的 Euler ステップで離散化される：
\eq{
	\itt{X} \eqdef \it{X} - \tau \nabla \Ff(\it{X}).
}
図~\ref{fig-flow-lagr}は、線形項とエントロピー項の和の汎関数に対するこのような離散化された陽的発展の例を示しており、Fokker--Planck 方程式の離散化版に対応する。
この特定の線形 Fokker--Planck 方程式の場合、確率的 PDE 手法に頼ることも可能であり、ガウスドリフトを持つ単一のランダム粒子を発展させることで数値的に近似できる。これらのスキーム（いわゆる Langevin Monte Carlo）の定常分布への収束は、Wasserstein 距離の観点で定量化できる。例えば~\citep{dalalyan2017user}を参照されたい。
関数 $\Ff$ が滑らかでない場合、~\eqref{eq-implicit-euclidean}と同様に陰的 Euler ステップで離散化する必要がある。\ie
\eq{
	\itt{X} \eqdef \Prox_{\tau \Ff}^{\norm{\cdot}}(\it{X}) \eqdef \uargmin{Z \in \X^n} \frac{1}{2}\norm{Z-\it{X}}^2 + \tau \Ff(Z).
}
を考える。
線形関数 $F(\al) = \int_X V(x) \d\al(x)$ の最も単純な場合、流れは各粒子 $x_i(t)$ に対して独立に作用し、関数 $V$ に対する通常のユークリッド流 $x_i'(t)=-\nabla V(x_i(t))$ に対応する（そして流れの積分曲線に沿った密度の移流 PDE である）。
\end{rem1}

\newcommand{\FigLagrFlow}[1]{\includegraphics[width=.19\linewidth]{lagrangian-flows/lagrange-flow-#1}}


\begin{figure}[h!]
\centering
\begin{tabular}{@{}c@{}c@{}c@{}c@{}c@{}}
\FigLagrFlow{1} &
\FigLagrFlow{3} &
\FigLagrFlow{5} &
\FigLagrFlow{7} &
\FigLagrFlow{10} \\
$t=0$ & $t=0.2$  & $t=0.4$ & $t=0.6$ & $t=0.8$
\end{tabular}
\caption{\label{fig-flow-lagr}
関数 $F(\al) = \int V \d\al - H(\al)$（$V(x)=\norm{x}^2$）に対する Lagrange 的離散化を用いた勾配流発展の例。
エントロピーは~\eqref{eq-entropy-discretized}を用いて離散化される。
極限の定常分布はガウス分布である。
}
\end{figure}


%%%%%%%%%%%%%%%%%%
\begin{rem2}{測地的凸性}
勾配流に関連する重要な概念は、Wasserstein-2 幾何学に関する汎関数 $F$ の凸性、\ie Wasserstein 測地線（注意~\ref{rem-displacement}に示される変位補間）に沿った $F$ の凸性である。このような関数に対する Wasserstein 勾配流は（連続時間で）存在し、一意であり、$\tau \rightarrow 0$ のときの離散ステッピング~\eqref{eq-grad-flow-stepping-cont}の極限である。$t \rightarrow +\infty$ のとき固定された定常分布に収束する。
エントロピーは測地的凸関数の典型的な例であり、$V : \X \rightarrow \RR$, $W : \X \times \X \rightarrow \RR$ が凸関数であるような線形関数 $F(\al) = \int_\X V(x) \d\al(x)$ および二次相互作用関数 $F(\al) = \int_{\X \times \X} W(x,y) \d\al(x) \d\al(y)$ も同様である。
線形関数は古典的な意味で凸であるが、二次相互作用関数はそうとは限らないことに注意する。
典型的な例は $W(x,y)=\norm{x-y}^2$ であり、これは負半正定値カーネル（定義~\ref{def-negativedefinitekernel}を参照）であるため、$F(\al)$ は通常の意味で凹な関数に対応する（測地的には凸である）。
\citet{mccann1997convexity}の重要な結果は、$\X=\RR^\dim$ 上で $F(\al) = \int_{\RR^\dim} \phi(\density{\al}(x)) \d x$ の形の一般化「エントロピー」関数は、$\phi$ が凸で、$\phi(0)=0$、$\phi(t)/t \rightarrow +\infty$（$t \rightarrow +\infty$ のとき）、かつ $s \mapsto s^{\dim} \phi(s^{-\dim})$ が凸減少であれば、測地的凸であるというものである。
\end{rem2}
%%%%%%%%%%%%%%%%%%%

結果として得られる離散化された流れの数値的解法に関する重要な文献があり、ここではいくつかの代表的な文献のみを挙げる。1次元問題に対しては、OT が逆累積分布関数の二次汎関数であること（注意~\ref{rem-1d-ot-generic}を参照）を利用した非常に精密なソルバが開発されている：\citet{kinderlehrer1999approximation,blanchet2008convergence,agueh2013one,Matthes1D,blanchet2012optimal}。高次元では、有限要素法および有限体積法スキーム~\citet{CarrilloFiniteVolume,burger2010mixed}で取り組むことができる。代替のソルバは Lagrange 的スキーム（\ie 粒子系）を用いて得られる：\citet{carrillo2009numerical,JDB-JKO,Westdickenberg2010}。
もう1つの方向は、連続の対応物のいくつかの性質を維持する離散流（典型的には離散格子やグラフ上）を探すことである。\citet{MielkeCVPDE,ErbarDCDS,ChowHuangLiZhou2012,Maas2011}を参照されたい。


エントロピー正則化に基づく Euler 的離散化問題~\eqref{eq-explicit-euclidean}を解くための近似的アプローチは、\citet{2015-Peyre-siims}で最初に提案され、\citet{2016-chizat-sinkhorn}で洗練され、\citet{2017-carlier-SIMA}で理論的に解析された。
エントロピー正則化を用いると、問題~\eqref{eq-grad-flow-discr}は $G = \iota_{\it{\a}}$ と設定し $F$ を $\tau F$ で置き換えたときの~\eqref{eq-generalized-ot-regul}の形を持つ。したがって、\citet{2015-Peyre-siims}で最初に提案されたように、反復~\eqref{eq-gen-sinkh}を用いて $\itt{\a}$ を近似できる。$\epsilon \rightarrow 0$ のときのこのスキームの収束は~\citet{2017-carlier-SIMA}で証明されている。
図~\ref{fig-jko}はこの方法で計算された発展の例を示す。
勾配流の機械学習への興味深い応用は、ある動力学的密度モデルを最もよくモデル化する基礎関数 $F$ を学習することである。この学習は、エントロピー正則化 transport と自動微分（注意~\ref{rem-auto-diff}を参照）を用いて滑らかな非凸最適化を解くことで達成できる。~\citet{hashimoto2016learning}を参照されたい。

時間と空間の両方で離散化された勾配流の収束を解析することは一般に困難である。距離を定義する線形計画問題の多面体的性質により、過度に小さなステップサイズの使用は「ロッキング」現象（分布が動かなくなり発展しない）につながるため、ステップサイズは小さすぎてはならない（~\citep{maury201713}で議論されている）。1次元における勾配流の離散化手法の収束解析については~\citep{Matthes1D,matthes2017convergent}を参照されたい。


\begin{figure}[h!]
\centering
\begin{tabular}{@{}c@{\hspace{1mm}}c@{\hspace{1mm}}c@{\hspace{1mm}}c@{\hspace{1mm}}c@{}}
\includegraphics[width=.19\linewidth]{jko/congestion/holes-potential}&
\includegraphics[width=.19\linewidth]{jko/congestion/holes-kappa10-1}&
\includegraphics[width=.19\linewidth]{jko/congestion/holes-kappa10-2}&
\includegraphics[width=.19\linewidth]{jko/congestion/holes-kappa10-10}&
\includegraphics[width=.19\linewidth]{jko/congestion/holes-kappa10-20}\\
\includegraphics[width=.19\linewidth]{jko/surfaces/cosw}&
\includegraphics[width=.19\linewidth]{jko/surfaces/evol-1}&
\includegraphics[width=.19\linewidth]{jko/surfaces/evol-2}&
\includegraphics[width=.19\linewidth]{jko/surfaces/evol-3}&
\includegraphics[width=.19\linewidth]{jko/surfaces/evol-4}\\
$\cos(w)$ & $t=0$ & $t=5$ & $t=10$ & $t=20$
\end{tabular}
\caption{\label{fig-jko}
ドリフト $V$ と混雑項を持つ勾配流発展の例（\citet{2015-Peyre-siims}より）。$F(\al) = \int_\X V(x) \d \al(x) + \iota_{\leq \kappa}(\density{\al})$。
}
\end{figure}



第~\ref{sec-unbalanced}節で詳述される非平衡 optimal transport 距離に対しても勾配流を計算することが可能である。これにより質量の生成または消滅を許容する発展が得られ、多くの物理、生物、化学現象をモデル化するために重要である。非平衡勾配流の例として、細胞成長のための著名な Hele-Shaw モデル~\citep{PerthameTumor}があり、理論的には~\citep{gallouet2017jko,marino2017tumor}で研究されている。このような非平衡勾配流も一般化 Sinkhorn アルゴリズム~\citep{2016-chizat-sinkhorn}を用いて近似できる。


%%%%%%%%%%%%%%%%%%%%%%%%%%%%%%%%%%%%%
%%%%%%%%%%%%%%%%%%%%%%%%%%%%%%%%%%%%%
%%%%%%%%%%%%%%%%%%%%%%%%%%%%%%%%%%%%%
%%%%%%%%%%%%%%%%%%%%%%%%%%%%%%%%%%%%%
\section{最小 Kantorovich 推定量}
\label{sec-mke}

未知の分布からの離散サンプル $(x_i)_{i=1}^n \subset \Xx$ が与えられたとき、目標はパラメトリックモデル $\th \mapsto \al_\th \in \Mm(\Xx)$ を観測された経験入力測度 $\be$ にフィッティングすることである：
\eql{\label{eq-density-fitting}
	\umin{\th \in \Theta} \Ll(\al_\th,\be)
	\qwhereq
	\be = \frac{1}{n} \sum_i \de_{x_i},
}
ここで $\Ll$ は離散分布と「連続的な」（任意の）分布との間のある「損失」関数である（図~\ref{fig-density-fitting}を参照）。

$\al_\th$ が Lebesgue 測度（またはその他の固定された参照測度）に関する密度 $\density{\th} \eqdef \density{\al_\th}$ を持つ場合、最尤推定量（MLE）は以下を解くことで得られる：
\eq{
	\umin{\th} \Ll_{\text{MLE}}(\al_\th,\be) \eqdef -\sum_i \log(\density{\th}(x_i)).
}
これは、$x_i$ がある $\bar\be$ のi.i.d.サンプルであると仮定すると、
\eq{
	\Ll_{\text{MLE}}(\al,\be) \overset{n \rightarrow +\infty}{\longrightarrow} \KL(\al|\bar\be)
}
となるため、Kullback--Leibler 損失の経験的対応物の使用に対応する。

\begin{figure}[h!]
\centering
\includegraphics[width=.4\linewidth]{density-fitting/schematic-fitting}
\caption{\label{fig-density-fitting}
密度フィッティング問題~\ref{eq-density-fitting}の模式図。
}
\end{figure}


\newcommand{\fPF}{h}

この MLE アプローチは多くの場合に最適な推定手続きにつながることが知られている（例えば~\citet{owen2001empirical}を参照）。しかし、特異な分布を推定する場合、典型的には $\al_\th$ が密度を持たないとき（$\Ll_{\text{MLE}}(\al_\th,\be) = +\infty$）、または $(x_i)_i$ がある特異な $\bar\be$ からのサンプルであるとき（$\KL(\al_\th|\bar\be)$ が有限であるためには $\al_\th$ が $\be$ と同じ台を共有する必要があるが、この台は通常未知である）には機能しない。もう1つの問題は、いくつかの実用上の関心がある場合において、密度 $\density{\th}$ にアクセスできない（または計算が困難すぎる）ことである。

両方の問題（特異性と未知の密度）が生じる典型的な設定は、いわゆる生成モデルであり、パラメトリック測度が固定された参照測度 $\zeta \in \Mm(\Zz)$ の押し出しとして書かれる：
\eq{
	\al_\th = \fPF_{\th,\sharp} \zeta \qwhereq \fPF_\th : \Zz \rightarrow \Xx,
}
ここで押し出し作用素は定義~\ref{defn-pushfwd}で導入される。空間 $\Zz$ は通常低次元であり、$\al_\th$ の台は低次元の「多様体」に沿って局在化し、結果として得られる密度は高度に特異的である（Lebesgue 測度に関する密度を持たない）。
さらに、この密度の計算は通常困難であるが、$\al_\th$ からi.i.d.サンプルを生成することは $x_i=\fPF_\th(z_i)$ を計算することで達成される。ここで $(z_i)_i$ は $\zeta$ からのi.i.d.サンプルである。

このような困難なシナリオに対処するためには、MLE 汎関数 $\Ll_{\text{MLE}}$ の代わりに弱い距離を用い、双対形式として
\eql{\label{eq-dual-loss}
	\Ll(\al,\be) \eqdef
	\umax{(\f,\g) \in \Cc(\X)^2}
	\enscond{ \int_{\X} \f(x) \d\al(x) + \int_{\X} \g(x) \d\be(x) }{ (\f,\g) \in \Potentials}
}
と書く必要がある。
第~\ref{sec-dual-norms}節で示される双対ノルムは
\eq{
	\Potentials = \enscond{(\f,-\f)}{\f \in B},
}
を課すことに対応し、optimal transport~\eqref{eq-dual-generic}は~\eqref{eq-dfn-pot-dual}で定義される $\Potentials = \Potentials(\c)$ を設定する。

固定された $\th$ に対して、~\eqref{eq-density-fitting}で最小化すべきエネルギーをこのような損失関数を用いて評価することは半離散 optimal transport を解くことに対応し、これは第~\ref{c-algo-semidiscr}章の焦点である。$\th$ に関するエネルギーの最小化はさらに困難であり、典型的には高度に非凸である。

$\Ee(\th) = \Ll(\al_\th,\be)$ を評価する際の~\eqref{eq-dual-loss}の解を $\f_\th$ と記すと、劣勾配が以下の公式で得られる：
\eql{\label{eq-grad-wass-loss-dual}
	\nabla \Ee(\th) = \int_\Xx [\partial \fPF_\th(x)]^\top \nabla \f_\th(x) \d\al_\th(x),
}
ここで $\partial \fPF_\th(x) \in \RR^{\text{dim}(\Theta) \times \dim}$ は $\th \in \RR^{\text{dim}(\Theta)} \mapsto \fPF_\th(x)$ の（$\th$ に関する）微分であり、$\nabla \f_\th(x)$ は $\f_\th$ の（$x$ に関する）勾配である。
この公式は数値的に使用が困難である。第一に、半離散問題の解である連続関数 $\f_\th$ を最初に計算する必要がある。第~\ref{sec-entropy-ot-mmd}節で示されるように、OT 損失に対してこれは確率的最適化を用いて達成できるが、高次元ではほとんど適用できない。もう1つの選択肢は、このポテンシャルにパラメトリック形式を課すことであり、例えば RKHS での展開（\citet{genevay2016stochastic}）やディープネットワーク近似（\citep{WassersteinGAN}）がある。しかし、これは理論的にはまだ解析されていない重要な近似誤差につながる。
最後の問題は、双対ポテンシャル $\f_\th$ の勾配 $\nabla \f_\th$ の計算を必要とするため、数値的に不安定であることである。

OT 損失に対しては、以下の等価な問題に対する主最適カップリングを計算すると、別の勾配公式が得られる：
\eql{\label{eq-ot-couplings-fitting}
	\MK_\c(\al_\th,\be) = \umin{\ga \in \Mm(\Zz \times \Xx)}
		\enscond{
			\int_{\Zz \times \Xx} c(\fPF_\th(z),x) \d\ga(z,x)
		}{
			\gamma \in \Couplings(\zeta,\be)
		}.
}
ここで考える半離散の場合、最小化すべき目的関数は実際に
\eql{\label{eq-ot-couplings-fitting-bis}
	\umin{(\ga_i)_{i=1}^n} \sum_{i=1}^n \int_\Zz c(\fPF_\th(z),x_i) \d \ga_i(z)
	\qwhereq
	\sum_{i=1}^n \ga_i = \zeta, \quad
	\int_\Zz \d\ga_i(z)  = \frac{1}{n},
}
と分解でき、各 $\ga_i \in \Mm_+^1(\Zz)$ である。
~\eqref{eq-ot-couplings-fitting-bis}を解く最適な $(\gamma_{\th,i})_i$ が得られると、$\Ee(\th)$ の勾配は
\eq{
	\nabla \Ee(\th) = \sum_{i=1}^n \int_\Zz [\partial \fPF_\th(z)]^\top \nabla_1 c(\fPF_\th(z),x_i) \d \ga_i(z),
}
と計算される。ここで $\nabla_1 c(x,y) \in \RR^\dim$ は $x \mapsto c(x,y)$ の勾配である。~\eqref{eq-grad-wass-loss-dual}とは異なり、この公式は双対 OT 問題の解であるポテンシャルの勾配の計算を含まないことに注意する。

$\Ll=\MK_\c$ を用いて得られる推定量のクラスは、しばしば「最小 Kantorovich 推定量」と呼ばれ、~\citep{bassetti2006minimum}で最初に導入された。~\citep{CanasRosasco}も参照されたい。これは制限付き Boltzmann マシンを訓練するために~\citep{CuturiBoltzman}により生成モデルの文脈で用いられ、~\citep{bernton2017inference}では近似的 Bayes 計算と組み合わせて用いられている。ディープネットワークを用いたこれらの計算の近似は、GAN~\citep{WassersteinGAN}と VAE~\citep{Bousquet2017}の両方のための深層生成モデルの訓練に用いられている。~\citep{2017-Genevay-AutoDiff,2017-Genevay-gan-vae,salimans2018improving}も参照されたい。パラメトリックモデルフィッティングへの Sinkhorn ダイバージェンスの使用は、形状マッチングおよび位置合わせに日常的に用いられている。~\citep{gold1998new,chui2000new,myronenko2010point,2017-feydy-miccai}を参照されたい。


\begin{rem1}{距離学習と転移学習}
	機械学習への応用において、OT 関連手法の成功はデータの幾何学を捉える適応されたコスト $c(x,y)$ の選択に大きく依存することを強調しておく。
	さまざまな種類のデータをユークリッド空間に埋め込むことは可能であるが（例えば単語埋め込みについては~\citep{mikolov2013efficient}を参照）、多くの場合、ある種の距離の適応や最適化が必要である。
	教師付きタスクのための距離学習は古典的な問題であり（例えば~\citep{MAL-019,weinberger2009distance}を参照）、学習パイプラインで何らかの OT 距離が用いられるときの ground 距離 $c(x,y)$ の学習に拡張されている~\citep{CuturiGroundMetric2014}（~\citealt{ZenICPR14,WangECCV12OLD,huang2016supervised}も参照）。
	また、最適カップリング $\P$ の観測からコスト行列を学習するという関連する逆問題も挙げておく。これは低ランク事前分布を用いて正則化できる~\citep{dupuy2016estimating}。
	関連する問題は転移学習~\citep{pan2010survey}とドメイン適応~\citep{glorot2011domain}であり、ある訓練済み機械学習パイプラインを新しいデータセットに適応させたい場合である。この問題は OT 技法を用いてモデル化し解くことができる。~\citep{courty2017optimal,courty2017joint}を参照されたい。
\end{rem1}
